\notechapter{N-20220808}{Set Counting, Modular Patterns, and Mathematical Induction}

\section{A set-counting problem}

The first page studies finite sets \(A_1,A_2,\ldots,A_{32}\) and asks how many subsets can satisfy certain intersection conditions.  The handwritten solution builds a toy example in which each set contains a fixed number of elements and one carefully controls pairwise intersections.

The general strategy is inclusion--exclusion.  If every set has \(m\) elements and pairwise overlaps have controlled size, then the size of the union is not simply the sum of sizes.  One must subtract overlaps and add back higher overlaps.  The note's construction eventually counts a large number of possible sets, obtaining a count of the form
\[
  31+28\cdot30=871.
\]
The exact model is less important than the idea: when many sets overlap, one should count by layers, not by naive addition.

\section{A modular arithmetic construction}

Another problem constructs two sets:
\[
  A=\{1,4,6,7,9,\ldots,1998\},
\]
and a reversed or complementary set near \(2000\).  The note observes that the pattern is governed by residues modulo \(11\).  Since
\[
  2000=11\cdot181+9,
\]
one can write elements as
\[
  11t+r.
\]
The allowed residues are something like
\[
  r\in\{1,4,6,7,9\}.
\]

The real insight is that forbidden differences correspond to forbidden residue differences.  Instead of checking hundreds of numbers, one checks a small table modulo \(11\).  This is one of the most common olympiad uses of modular arithmetic.

\section{Why not every construction is finished}

The page explicitly says ``problem reserved'' in a few places.  That is an honest part of mathematical notes.  The important thing is to record what was tried and why it seemed plausible.  Here the attempt is to remove residues that conflict pairwise and leave a maximal collection.  Whether the construction is optimal requires a separate upper-bound argument.

This distinction is crucial: constructing an example proves a lower bound, but proving maximality requires an upper bound.

\section{Sequence exercises}

Later pages return to sequence problems.  Examples include arithmetic and geometric sequences, partial sums, and recurrence-like identities.  One problem has
\[
  a_n=a_1+(n-1)d
\]
and conditions such as
\[
  a_1+a_2+\cdots+a_n=S_n.
\]
The correct habit is to translate everything into \(a_1,d,n\), solve, and then check the requested value.

Another problem uses a geometric progression and asks for a common ratio or a sum.  The formula
\[
  S_n=a_1\frac{q^n-1}{q-1}
\]
again appears, with the warning that \(q=1\) is a separate case.

\section{Mathematical induction}

The most substantial later pages discuss induction.  The note distinguishes several forms:
\begin{enumerate}[(i)]
  \item ordinary induction: prove \(P(1)\), then prove \(P(n)\Rightarrow P(n+1)\);
  \item strong induction: prove \(P(1),\ldots,P(n)\Rightarrow P(n+1)\);
  \item induction with a shifted start: begin at \(n=m\) instead of \(1\).

\end{enumerate}

The proof skeleton is:
\begin{enumerate}[(i)]
  \item base case;
  \item induction hypothesis;
  \item induction step;
  \item conclusion.

\end{enumerate}
The note emphasizes that the induction hypothesis must be used precisely.  One cannot simply assume the desired statement for \(n+1\).

\section{Examples of induction}

One exercise proves a formula for a sum, such as
\[
  1+2+\cdots+n=\frac{n(n+1)}2.
\]
The induction step is
\[
  1+\cdots+n+(n+1)
  =\frac{n(n+1)}2+(n+1)
  =\frac{(n+1)(n+2)}2.
\]

Another exercise proves a divisibility statement.  For such problems, the usual technique is to write the \(n+1\) case in terms of the \(n\) case plus a visibly divisible remainder.

The page also contains a problem of the form
\[
  4^{2k+1}+3^{k+1}
\]
or similar, where the induction step uses modular or factorization identities.  The guiding rule is: when proving divisibility by induction, preserve a factor rather than expand everything.


\section{Induction structure}

A proof by induction has two parts: the base case and the inductive step.  The inductive step is not simply checking the next case; it is proving
\[
  P(n)\Rightarrow P(n+1)
\]
for arbitrary \(n\).  This is why writing the induction hypothesis explicitly matters.

Strong induction allows the step to use all previous cases:
\[
  P(1),P(2),\ldots,P(n)\Rightarrow P(n+1).
\]
It is useful when the object of size \(n+1\) breaks into smaller pieces of varying sizes.

\section{Modular patterns}

When a problem depends on remainders, the natural universe is \(\mathbb Z/m\mathbb Z\).  Instead of listing many cases, one should identify the period.  If a sequence satisfies
\[
  a_{n+m}=a_n,
\]
then all behavior is determined by the first \(m\) terms.  This is the organizing idea behind many modular counting exercises.

\section{A note on scope}

This note is about finite structure.  Set problems use inclusion--exclusion and modular residues; sequence problems use closed forms; induction proves that a pattern persists for all \(n\).  The conceptual chain is:
\[
  \text{example}\quad\Rightarrow\quad\text{pattern}\quad\Rightarrow\quad\text{proof by induction or upper bound}.
\]
A mature solution must know which stage it is in.  Guessing a pattern is not proving it; constructing an example is not proving optimality; verifying a few cases is not induction.


\paragraph{Expanded synthesis.}
The number-theoretic content of this chapter is best understood as arithmetic seen through algebraic structure. The earlier sections (`A set-counting problem`, `A modular arithmetic construction`, `Why not every construction is finished`, `Sequence exercises`, `Mathematical induction`) compute divisibility, congruences, residues, prime factorizations, or finite-field examples. The deeper language is ideals, quotient rings, unit groups, valuations, and field extensions.

Divisibility in $\mathbb Z$ is containment of principal ideals; congruence modulo $n$ is equality inside the quotient ring $\mathbb Z/n\mathbb Z$; invertibility modulo $n$ means being a unit; the Chinese remainder theorem is a product decomposition of quotient rings. Prime factorization supplies coordinates through valuations:
\[
  n=\prod_p p^{v_p(n)}.
\]
Gcd and lcm are coordinatewise minimum and maximum in these valuation coordinates.

This algebraic reading explains why elementary tricks scale upward. Euclidean algorithms become the theory of Euclidean domains and PIDs; divisor sums become Dirichlet convolution and Euler products; roots of unity become cyclotomic fields and Galois actions; finite polynomial arithmetic becomes finite-field theory. The original computations should therefore be kept as the algorithmic core, while the advanced language identifies the structure that makes the algorithms work.
